\section{Introduction}
\label{sec:intro}

Which convex polytope with a given number of vertices, inscribed in a sphere,
has the largest volume? Fejes T\'{o}th raised this question in three
dimensions \cite{fejes-toth-1964}, and it remains open there for all but a few
vertex counts. Berman and Hanes found a local necessary condition for
optimality and used it to determine the optimal polyhedra with at most eight
vertices \cite{berman-hanes-1970}; Hoehner and Ledford recently settled the
case of nine vertices in a preprint \cite{hoehner-ledford-2026}. Beyond nine
vertices only numerical results are known: Mutoh reported a computer search
for up to $30$ vertices \cite{mutoh-2003}, and Hardin, Sloane and Smith
tabulate putatively optimal configurations for up to $130$
\cite{hardin-sloane-smith}.

In higher dimensions even less is known. In dimension $d$, Horv\'{a}th and
L\'{a}ngi determined the optimal polytopes with $d+2$ vertices, and, when $d$
is odd, with $d+3$ vertices \cite{horvath-langi-2016}. When $d$ is even, their
argument for $d+3$ vertices covers only polytopes that are not cyclic, and
they asked whether an optimal polytope can be cyclic
\cite[Problem~2]{horvath-langi-2016}; Horv\'{a}th restates the question in
\cite[Problem~2.5]{horvath-2016-survey}.

We answer this question in the first open case, $d=4$, where the polytopes
have seven vertices. We show that every cyclic polytope with seven vertices
inscribed in the unit sphere of $\mathbb{R}^4$ has volume at most $2/5$
(Theorem~\ref{thm:main}). A non-cyclic polytope attains the larger volume
$\sqrt{3}/4$, so no optimal polytope is cyclic (Corollary~\ref{cor:problem2}).
Combined with the result of Horv\'{a}th and L\'{a}ngi on non-cyclic
polytopes, this shows that the maximal volume is $\sqrt{3}/4$
(Corollary~\ref{cor:seven}).

We then settle the next case in four dimensions, eight vertices. Here the
natural candidate is the regular cross-polytope, the four-dimensional analogue
of the regular octahedron, which is optimal among polyhedra with six vertices
\cite{berman-hanes-1970}. Horv\'{a}th and L\'{a}ngi showed that in every
dimension the regular cross-polytope is optimal among polytopes with the
combinatorics of a cross-polytope \cite{horvath-langi-2016}, but this leaves
open whether a polytope of another combinatorial type does better. Since an
optimal polytope has only simplices as facets \cite{horvath-langi-2016}, it
suffices to consider simplicial polytopes. There are $37$ combinatorial types
of simplicial polytopes with eight vertices in four dimensions
\cite{grunbaum-sreedharan-1967}, and Firsching realised each of them with
rational vertices on the sphere \cite{firsching-2017}. We show that each of
the $36$ types other than the cross-polytope has volume at most $13/20$
(Theorem~\ref{thm:eight}), which is less than the volume $2/3$ of the regular
cross-polytope. Hence the regular cross-polytope is the optimal polytope with
eight vertices (Corollary~\ref{cor:eight}).

Both results rest on the same method. For a fixed combinatorial type we write
the volume as a polynomial in the coordinates of the vertices and bound it by
a sum-of-squares certificate \cite{parrilo-2003}, an algebraic identity that a
computer checks in exact rational arithmetic. The bound holds for every
configuration of points on the sphere, whether or not the points are in convex
position, so convexity never has to be verified.

Section~\ref{sec:prelim} collects the background, Section~\ref{sec:method}
describes the certificates and their verification, and
Sections~\ref{sec:results} and~\ref{sec:eight} prove the results for seven and
eight vertices.
